\section{Open questions}

The following five open questions came out of the replication and are not
answered by the paper as written. Each is stated with an evidentiary basis
(what in the paper or in this replication motivates it) and concrete next
steps that could be executed within our free-endpoint / classical-statevector
infrastructure.

\begin{enumerate}
    \item \textbf{Multi-amplitude extension.} Does MLE-QAE extend cleanly to
    multi-amplitude / multi-variable estimation, where the good-subspace is
    not 2-D but a superposition over several tagged outcomes?
    \emph{Basis:} Paper is restricted to a single scalar $a\in[0,1]$; real
    Monte Carlo pipelines often need joint estimation of several amplitudes.
    \emph{Next steps:} derive the $K$-outcome joint likelihood; simulate a
    3-amplitude case at statevector level; benchmark joint-MLE against $K$
    independent runs.

    \item \textbf{Interaction with error mitigation.} How does MLE-QAE
    integrate with ZNE / PEC / virtual distillation on NISQ hardware, and
    does mitigation preserve near-Heisenberg scaling or push it back toward
    $N^{-1/2}$?
    \emph{Basis:} paper is noiseless; depolarizing noise on $Q^{m_k}$ damps
    the signal exponentially in $m$, which is exactly the regime where EIS
    gains.
    \emph{Next steps:} add depolarizing noise to the good-subspace simulator;
    overlay a ZNE-style sweep; refit MLE and check whether slope stays below
    $-0.75$.

    \item \textbf{End-to-end vs.\ classical QMC on option pricing.} For the
    paper's Sec.\ 4 headline application (Monte Carlo option pricing), what
    is the empirical MSE-vs-total-queries curve for MLE-QAE against a fully
    classical QMC baseline at industrially-relevant $10^{-4}$ relative
    precision?
    \emph{Basis:} paper reports CNOT counts only, not head-to-head MSE-vs-
    queries against a classical baseline.
    \emph{Next steps:} wire the MLE-QAE post-processor to a small European-
    call payoff simulator; run both MLE-QAE and classical MC to matched MSE;
    report the crossover precision.

    \item \textbf{Robustness of MLE to model misspecification.} How robust is
    the MLE post-processing when the assumed $\sin^2((2m+1)\theta)$ likelihood
    is off due to readout-confusion, coherent over-rotations, or mid-circuit
    leakage?
    \emph{Basis:} MLE is only optimal when the likelihood model matches
    reality; hardware effects break this.
    \emph{Next steps:} perturb the forward model with a readout-confusion
    matrix and per-Q coherent over-rotation; measure bias and RMSE/CRB of the
    (unaware) MLE; test a data-augmented MLE that fits $\theta_a$ plus a
    readout-bias nuisance parameter.

    \item \textbf{Adaptive-schedule MLE-QAE.} Does an adaptive Bayesian
    scheduler for $\{m_k\}$ beat the fixed EIS schedule at matched query
    budget, especially in the low-shot regime where the finite-$N_{\rm shot}$
    bias currently makes the empirical slope shallower than the theoretical
    $-1$?
    \emph{Basis:} our replicated EIS slope is $-0.866$ vs.\ the paper's
    $-0.95$; the paper attributes the gap to finite-shot MLE bias, but does
    not test whether adaptive scheduling closes it.
    \emph{Next steps:} implement a Bayesian MLE-QAE loop that picks
    $m_{k+1}$ to maximize expected Fisher-information gain under the current
    posterior; compare RMSE-vs-$N_q$ against EIS at $N_{\rm shot}=100,1000$;
    report query-budget savings at fixed $10^{-5}$ relative precision.
\end{enumerate}
